\section{Conclusion}
\label{sec:conclusion}

We present the first controlled comparison of attractor formation across neural architecture families in the language domain. By training parameter-matched Transformer, GRU, and LSTM models on identical data with identical tokenization, we show that architecture alone produces measurably different default response patterns: matched-capability Transformer and GRU models agree on their top-1 prediction only 25\% of the time, with a mean Jensen-Shannon divergence of 0.34. At the representation level, Transformers create a universal attractor basin---a tight, undifferentiated representation space where all concepts converge---that offers a mechanistic explanation for persona collapse. Recurrent architectures, constrained by finite hidden state, preserve better concept separation.

Our findings suggest that the homogeneity of default behaviors in frontier LLMs is partly a consequence of the Transformer's inductive bias, not solely of shared training data or alignment procedures. This has practical implications: architectural diversity in deployed model ensembles may reduce persona collapse and associated monoculture risks.

Several promising directions remain. Scaling this analysis to modern recurrent architectures (Mamba, RWKV) at frontier scale would test whether the architecture--attractor relationship holds in the regime where it matters most. Applying RLHF to matched architectures would isolate alignment's contribution to persona collapse. Finally, layer-wise attractor analysis---extending the framework of \citet{chytas2025concept} to recurrent networks---could reveal \emph{where} in the network architecture-specific defaults emerge, informing targeted interventions.
